\bd[Programming]
\textbf{Programming} is the process of designing and building an executable computer program to accomplish a specific
computing result or to perform a particular task.
\ed

Programming involves tasks such as analysis, generating algorithms, profiling algorithms' accuracy and resource
consumption, and the implementation of algorithms in a chosen programming language (commonly referred to as coding).
The source code of a program is written in one or more languages that are intelligible to programmers, rather than
machine code, which is directly executed by the central processing unit.

\bd[Source Code]
\textbf{Source code} is any collection of code, with or without comments, written using a human-readable programming
language, usually as plain text.
\ed

The purpose of programming is to find a sequence of instructions that will automate the performance of a task (which
can be as complex as an operating system) on a computer, often for solving a given problem. Proficient programming
thus usually requires expertise in several subjects, including knowledge of the application domain,
specialized algorithms, and formal logic. Tasks accompanying and related to programming include testing, debugging,
source code maintenance, implementation of build systems, and management of derived artifacts, such as the machine
code of computer programs. \v

Programming languages are easy and efficient for programmers because they are closer to natural languages than the
machine code. Programmers use programming languages to write software.

\bd[Software]
\textbf{Software} is a collection of instructions that tell a computer how to work.
\ed

\bd[Software Development]
\textbf{Software development} is the process of conceiving, specifying, designing,programming, documenting, testing,
and bug fixing involved in creating and maintaining applications, frameworks, or other software components.
\ed

To build huge programs, programmers use a set of tools and practices. Taken together, these form the discipline of
``Software Engineering'', a term coined by engineer Margaret Hamilton, who helped NASA prevent serious problems on
the Apollo 11 mission to the moon.

\bd[Software Engineering]
\textbf{Software engineering} is the systematic application of engineering approaches to the development of software.
\ed

\bd[Software Engineer]
A \textbf{software engineer} is a person who applies the principles of software engineering to design, develop,
maintain, test, and evaluate computer software.
\ed

Software engineering combines engineering techniques with software development practices. \v

In this chapter we will be focusing on the fundamentals of programming and software engineering.

\section{Programming Basics}

Just like spoken languages, programming languages have statements.

\bd[Statement]
A \textbf{statement} is a syntactic unit of an imperative programming language that expresses some action to be carried
out.
\ed

A program written in such a language is formed by a sequence of one or more statements. A statement may have internal
components called ``expressions''.

\bd[Expression]
An \textbf{expression} is a syntactic entity in a programming language that may be evaluated to determine its value.
\ed

An expressions is a combination of one or more constants, variables, functions, and operators that the programming
language interprets and computes to return another value. This process, for mathematical expressions, is called
``evaluation''.

\bd[Constant]
A \textbf{constant} is a value that should not be altered by the program during normal execution, i.e.\ the value is
constant.
\ed

When associated with an identifier, a constant is said to be ``named'', although the terms ``constant'' and ``named
constant'' are often used interchangeably. Constants are useful for both programmers and compilers: For programmers
they are a form of self-documenting code and allow reasoning about correctness, while for compilers they allow
compile-time and run-time checks that verify that constancy assumptions are not violated, and allow or simplify some
compiler optimizations.

\bd[Operators]
\textbf{Operators} are constructs defined within programming languages which behave generally like functions, but
which differ syntactically or semantically.
\ed

Common simple examples include arithmetic like addition, comparison and greater than, and logical operations like AND.
More involved examples include assignments field access in a record or object, and the scope resolution operator.
Languages usually define a set of built-in operators, and in some cases allow users to add new meanings to existing
operators or even define new operators. \v

Constants and operators are contrasted with a variable, which is an identifier with a value that can be changed
during normal execution, i.e.\ , the value is variable.

\bd[Variable]
A \textbf{variable} is an abstract storage location paired with an associated symbolic name, which contains some
known or unknown quantity of information referred to as a value, or in simpler terms, a variable is a container for a
particular set of bits or type of data.
\ed

A variable can eventually be associated with or identified by a memory address. The variable name is the usual way to
reference the stored value, in addition to referring to the variable itself, depending on the context. This
separation of name and content allows the name to be used independently of the exact information it represents. The
identifier in computer source code can be bound to a value during run time, and the value of the variable may thus
change during the course of program execution.

\bd[Subroutine]
A \textbf{subroutine} is a sequence of program instructions that performs a specific task, packaged as a unit. This
unit can then be used in programs wherever that particular task should be performed.
\ed

Subroutines may be defined within programs, or separately in libraries that can be used by many programs. In
different programming languages, a subroutine may be called a routine, subprogram, function, method, or procedure.
Technically, these terms all have different definitions. \v

Software consists of thousands of smaller subroutines, each responsible for different features. In modern
programming, it's uncommon to see subroutines longer than around 100 lines of code, because by then, there's probably
something that should be pulled out and made into its own subroutine. Modularizing programs into subroutines not only
allows a single programmer to write an entire app, but also allows teams of people to work efficiently on even bigger
programs. Different programmers can work on different subroutines, and if everyone makes sure their code works
correctly, then when everything is put together, the whole program should work too! Modern programming languages come
with huge bundles of pre-written subroutines, called ``libraries''.

\bd[Library]
A \textbf{library} is a collection of resources used by computer programs, often for software development. These may
include configuration data, documentation, help data, message templates, pre-written code and subroutines, classes,
values or type specifications.
\ed

Libraries are written by expert coders, made efficient and rigorously tested,and then given to everyone. There are
libraries for almost everything, including networking, graphics, and sound -- topics we'll discuss in later chapters.
For now let's go back to the fundamentals. \v

The English language has syntax, and so do all programming languages.

\be
``A equals 5'' is a programming language statement. In this case, the statement says a variable named $A$ has the number
$5$ stored in it. This is called an assignment statement because we're assigning a value to a variable.
\ee

\bd[Syntax]
The set of rules that govern the structure and composition of statements in a language is called \textbf{syntax}.
\ed

To express more complex things, we need a series of statements.

\be
For example ``A is 5, B is ten, C equals A plus B''. This program tells the computer to set variable $A$ equal to
$5$, variable $B$ to 10, and finally to add $A$ and $B$ together, and put that result, which is 15, into variable $C$.
In general a computer doesn't care about the name of the variables, as long as variables are uniquely named. But
it's probably best practice to name them things that make sense in case someone else is trying to understand your
code. For best practices check Appendix~\ref{ch:clean_code}.
\ee

In general, a program starts at the first statement and runs down one at a time until it hits the end. One can alter
this behaviour by making use of the so called ``control flow statements''.

\bd[Control Flow Statement]
A \textbf{control flow statement} is a statement that results in a choice being made as to which of two or more paths
to follow.
\ed

There are several control flow statement types.

\bd[Conditional Expressions]
\textbf{Conditional expressions} are features of a programming language which perform different computations or
actions depending on whether a programmer-specified boolean condition evaluates to true or false.
\ed

The most used conditional expression is the ``If-Then-Else'' statements which are executed once, a conditional path
is chosen, and the program moves on.

\fig{if}{0.38}

To repeat some statements many times, we need to create a conditional loop.

\bd[Conditional Loops]
\textbf{Conditional loops} are a way for computer programs to repeat one or more various steps depending on
conditions set either by the programmer initially or real-time by the actual program.
\ed

One way is a ``while statement'', also called a ``while loop''. A ``while loop'' checks condition for truthfulness
before executing any of the code in the loop. If condition is initially false, the code inside the loop will never be
executed. \v

There's also the common ``for loop''. Instead of being a condition-controlled loop that can repeat forever until the
condition is false, a ``for loop'' is count-controlled; it repeats a specific number of times. \v

As we mentioned in the previous sections, breaking big programs into smaller functions allows many people to work
simultaneously. They don't have to worry about the whole thing, just the function they're working on. So,  if you're
tasked with writing a sort algorithm, you only need to make sure it sorts properly and efficiently. However, even
packing up code into functions isn't enough.

\be
Microsoft Office probably contains hundreds of thousands of functions. That's better than dealing with 40 million
lines of code, but it's still way too many ``things'' for one person or team to manage.
\ee

The solution is to package functions into hierarchies, pulling related code together into ``objects''.

\bd[Object]
An \textbf{object} is the unit of code that is eventually derived from the process. It can be a variable, a data
structure, a function, or a method.
\ed

\be
For example, car's software might have several functions related to cruise control, like setting speed, nudging speed
up or down, and stopping cruise control altogether. Since they're all related, we can wrap them up into a unified
cruise control object. \v

But, we don't have to stop there. Cruise control is just one part of the engine's software. There might also be sets
of functions that control spark plug ignition, fuel pumps, and the radiator. So we can create a ``parent'' Engine
object that contains all of these ``children'' objects. \v

In addition to children objects, the engine itself might have its own functions. You want to be able to stop and
start it, for example. It'll also have its own variables, like how many miles the car has traveled. In general,
objects can contain other objects, functions, and variables. And of course, the engine is just one part of a Car
object. There's also the transmission, wheels, doors, windows, and so on. \v

Now, as a programmer, if I want to set the cruise control, I navigate down the object hierarchy, from the outermost
objects to more and more deeply nested ones. Eventually, I reach the function I want to trigger: Car, then engine,
then cruise control, then set cruise speed to 55.
\ee

Programming languages often use something equivalent to the syntax described in the example. The idea of packing up
functional units into nested objects is called Object Oriented Programming.

\bd[Object Oriented Programming (OOP)]
\textbf{Object-oriented programming} (\textbf{OOP}) is a programming paradigm based on the concept of ``objects'',
which can contain data and code.
\ed

OOP is very similar to what we've been doing all long: hide complexity by encapsulating low-level details in
higher-order components. Before we packed up things like transistor circuits into higher-level boolean gates. Now,
we're doing the same thing with software. Yet again, it's a way to move up a new level of abstraction. \v

OOP encapsulates data in the form of fields, often known as ``attributes'' or ``properties'', and code in the form of
procedures, often known as ``methods''.

\bd[Attribute / Property]
An \textbf{attribute} or \textbf{property} is a specification that defines a property of an object.
\ed

For clarity, attributes should more correctly be considered metadata. An attribute is frequently and generally a
property of a property. However, in actual usage, the term attribute can and is often treated as equivalent to a
property depending on the technology being discussed. An attribute of an object usually consists of a name and a value.

\bd[Method]
A \textbf{method} is a function associated with an object.
\ed

Breaking up a big program, like a car's software, into functional units is perfect for teams. One team might be
responsible for the cruise control system, and a single programmer on that team tackles a handful of functions. This
is similar to how big, physical things are built, like skyscrapers. You'll have electricians running wires, plumbers
fitting pipes, welders welding, painters painting, and hundreds of other people teeming all over the hull. They work
together on different parts simultaneously, leveraging their different skills, until one day, you've got a whole
working building.

A method cain be either public or private.

\bd[Public Method]
\textbf{Public methods} are methods that are accessible both inside and outside the scope of your class. Any instance
of that class will have access to public methods and can invoke them.
\ed

\bd[Private Method]
\textbf{Private methods} are those methods which can't be accessed in other class except the class in which they are
declared. We can perform the functionality only within the class in which they are declared.
\ed

\be
Returning to our cruise control example: its code is going to have to make use of functions in other parts of the
engine's software to, you know, keep the car at constant speed. That code isn't part of the cruise control team's
responsibility. It's another team's code. Because the cruise control team didn't write that, they're going to need
good documentation about what each function in the code does, and a well-defined Application Programming Interface,
or API for short.
\ee

\bd[Documentation]
\textbf{Documentation} is any communicable material that is used to describe, explain or instruct regarding some
attributes of an object, system or procedure, such as its parts, assembly, installation, maintenance and use.
\ed

Documentation can be provided on paper, online, or on digital or analog media, such as audio tape or CDs. Examples
are user guides, white papers, online help, and quick-reference guides. Paper or hard-copy documentation has become
less common. Documentation is often distributed via websites, software products, and other online applications. \v

Documentation promotes code reuse. So, instead of having programmers constantly write the same things over and over,
they can track down someone else's code that does what they need. Then, thanks to documentation, they can put it to
work in their program, without ever having to read through the code.

\bd[Application Programming Interface (API)]
An \textbf{Application Programming Interface} (\textbf{API}) is a connection between computers or between computer
programs. It is a type of software interface, offering a service to other pieces of software.
\ed

\bd[API Specification]
An \textbf{API specification} is a document or standard that describes how to build or use an API. A computer system
that meets this standard is said to implement or expose an API.
\ed

The term API may refer either to the specification or to the implementation. \v

An API connects computers or pieces of software to each other. It is not intended to be used directly by a person
(the end user) other than a computer programmer who is incorporating it into software. An API is often made up of
different parts which act as tools or services that are available to the programmer. A program or a programmer that
uses one of these parts is said to call that portion of the API. The calls that make up the API are also known as
subroutines, methods, requests, or endpoints. An API specification defines these calls, meaning that it explains how
to use or implement them. \v

One purpose of APIs is to hide the internal details of how a system works, exposing only those parts a programmer
will find useful and keeping them consistent even if the internal details later change. An API may be custom-built
for a particular pair of systems, or it may be a shared standard allowing interoperability among many systems.

\be
For example, in the Ignition Control object, there might be functions to set the RPM of the engine, check the spark
plug voltage, as well as fire the individual spark plugs. Being able to set the motor's RPM is really useful. The
cruise control team is going to need to call that function, but they don't know much about how the ignition system
works. It's not a good idea to let them call functions that fire the individual spark plugs, or the engine might
explode! Maybe. The API allows the right people access to the right functions and data.
\ee

OOP languages allow the right people to have access to the right functions and data by letting you specify whether
methods are public or private. This ability to hide complexity, and selectively reveal it,  is the essence of OOP,
and it's a powerful and popular way to tackle building large and complex programs. Pretty much every piece of
software on your computer, or game running on your console was built using an OOP language, like C++, C\#, or
Objective-C. Other popular OOP languages you may have heard of are Python and Java. It's important to remember that
code, before being compiled, is just text. As I mentioned earlier, you could write code in Notepad or any old word
processor. Some people do. But generally, today's software developers use special-purpose applications for writing
programs, ones that integrate many useful tools for writing, organizing, compiling and testing code. Because they've
put everything you need in one place, they're called Integrated Development Environments, or IDEs for short.

\bd[Integrated Development Environment (IDE)]
An \textbf{Integrated Development Environment} (\textbf{IDE}) is a software application that provides comprehensive
facilities to computer programmers for software development.
\ed

An IDE normally consists of at least a source code editor for writing code (often with useful features like automatic
color-coding to improve readability, syntax error checking, and many more), build automation tools and a debugger.
Some IDEs, such as NetBeans and Eclipse, contain the necessary compiler, interpreter, or both; others, such as
SharpDevelop and Lazarus, do not. The boundary between an IDE and other parts of the broader software development
environment is not well-defined; sometimes a version control system or various tools to simplify the construction of
a graphical user interface (GUI) are integrated. Many modern IDEs also have a class browser, an object browser, and a
class hierarchy diagram for use in object-oriented software development. \v

IDE can take you back to the line of code where a bug happened, and often provide you additional information to help
you track down and fix the bug. This is important because most programmers spend 70 to 80\% of their time testing and
debugging, not writing new code. Good tools, contained in IDEs, can go a long way when it comes to helping
programmers prevent and find errors. \v

In addition to coding and debugging, another important part of a programmer's job is documenting their code. This can
be done in standalone files called ``README'' which tell other programmers to read that help file before diving in.

\bd[README]
A \textbf{README} file contains information about the other files in a directory or archive of computer software. A
form of documentation, it is usually a simple plain text file. The file's name is generally written in uppercase.
\ed

It can also happen right in the code itself with comments. These are specially-marked statements that the program
knows to ignore when the code is compiled. They exist only to help programmers figure out what's what in the source
code. Good documentation helps programmers when they revisit code they haven't seen for a while, but it's also crucial
for programmers who are totally new to it. \v

In addition to IDEs, another important piece of software that helps big teams work collaboratively on big coding
projects is called Version Control.

\bd[Version Control]
\textbf{Version control} (also known as revision control, source control, or source code management) is a class of
systems responsible for managing changes to computer programs, documents, large websites, or other collections of
information. Version control is a component of software configuration management.
\ed

A version control can be either centralized or distributed.

\bd[Centralized Version Control]
\textbf{Centralized version control} systems are based on the idea that there is a single, central copy of a project
on a server and programmers will commit their changes to this central copy.
\ed

The most popular centralized version control systems are Git and Mercurial

\bd[Distributed Version Control]
\textbf{Distributed version control} systems are based on the idea that the complete codebase — including its full
version history — is mirrored on every developer's computer.
\ed

The most popular distributed version control systems are Subversion, CVS, and Perforce.

\bd[Repository]
A \textbf{repository} is a data structure that stores metadata for a set of files or directory structure.
\ed

Depending on whether the version control system in use is distributed or centralized the whole set of information in
the repository may be duplicated on every user's system or may be maintained on a single server. Some of the metadata
that a repository contains includes, among other things a historical record of changes in the repository, a set of
commit objects. and a set of references to commit objects, called heads. \v

When a programmer wants to work on a piece of code, they can check it out, sort of like checking out a book out from
a library. Often, this can be done right in an IDE. Then, they can edit this code all they want on their personal
computer, adding new features and testing if they work. When the programmer is confident their changes are working
and there are no loose ends, they can check the code back into the repository, known as committing code, for everyone
else to use. While a piece of code is checked out, and presumably getting updated or modified, other programmers
leave it alone. This prevents weird conflicts and duplicated work. In this way, hundreds of programmers can be
simultaneously checking in and out different pieces of code, iteratively building up huge systems. \v

Critically, you don't want someone committing buggy code, because other people and teams may rely on it. Their code
could crash, creating confusion and lost time. The master version of the code, stored on the server, should always
compile without errors and run with minimal bugs. But sometimes bugs do creep in. Fortunately, source control
software keeps track of all changes, and if a bug is found, the whole code, or just a piece, can be rolled back to an
earlier, stable version. It also keeps track of who made each change, so coworkers can send helpful and encouraging
emails to the offending person. \v

Debugging goes hand in hand with writing code, and it's most often done by an individual or small team. The big
picture version of debugging is Quality Assurance testing, or QA.

\bd[Quality Assurance (QA)]
\textbf{Quality Assurance} (\textbf{QA}) is any systematic process of determining whether a product or service meets
specified requirements.
\ed

QA is where a team rigorously tests out a piece of software, attempting to create unforeseen conditions that might
trip it up. Basically, they elicit bugs. Getting all the wrinkles out is a huge effort, but vital in making sure the
software works as intended for as many users in as many situations as imaginable before it ships. \v

You've probably heard of beta software.

\bd[Beta Software]
\textbf{Beta software} refers to computer software that is undergoing testing and has not yet been officially
released.
\ed

Companies will sometimes release beta versions to the public to help them identify issues. It's essentially like
getting a free QA team.

% TODO: WIP
\section{Algorithms \& Data Structures}

\bd[Data]
\textbf{Data} are the quantities, characters, or symbols on which operations are performed by a computer, which may
be stored and transmitted in the form of electrical signals and recorded on magnetic, optical, or mechanical
recording media.
\ed

If data are arranged in a systematic way then they get a structure and become meaningful. These meaningful (or
processed) data are called ``information''.

\bd[Information]
\textbf{Information} is processed, organized and structured data.
\ed

Information provides context for data and enables decision-making process.

\be
For example, a single customer’s sale at a restaurant is data, this becomes information when the business is able to
identify the most popular or least popular dish.
\ee

A data structure is a systematic way to organize data so that they can be used efficiently.

\bd[Data Structure]
A \textbf{data structure} is a data organization, management,and storage format that enables efficient access and
modification.
\ed

More precisely, a data structure is a collection of data values, the relationships among them, and the functions or
operations that can be applied to the data, i.e.\ it is an algebraic structure about data. Data structures are
essential ingredients in creating fast and powerful algorithms, and they make code cleaner and easier to understand. \v

\bd[Abstract Data Type (ADT)]
An \textbf{abstract data type} (\textbf{ADT}) is a mathematical model for data types which is defined by its behavior
(semantics) from the point of view of a user of the data and more specifically in terms of possible values, possible
operations on data of this type, and the behavior of these operations.
\ed

ADTs contrast with data structures, which are concrete representations of data, and are the point of view of an
implementer, not a user. \v

Formally, an ADT may be defined as a ``class of objects whose logical behavior is defined by a set of values and a
set of operations'', which is analogous to an algebraic structure in mathematics. In practice, many common data types
are not ADTs, as the abstraction is not perfect. \v

\textbf{Data structures are used to implements ADTs.} There are multiple ways (i.e.\ data structures) to implement an
ADT. In simple words, ADTs tell us what is to be done and data structures tell us how to do it. In reality, different
implementations of ADTs (a.k.a different data structures) are compared for time and space efficiency. The one best
suited depends on the requirements of the project. \v

Data structures are divided into categories based on their characteristics. One basic division is between ``linear''
and ``non-linear'' data structures.

\bd[Linear Data Structures]
\textbf{Linear data structures} have data elements arranged in linear (sequential) order and each member element is
connected to its previous and next element, except from the first and the last.
\ed

\be
Examples of linear data structures are lists, queues, stacks and arrays.
\ee

\bd[Non-Linear Data Structures]
\textbf{Non-linear data structures} are those data structures in which data items are not arranged in a sequence.
\ed

\be
Examples of non-linear data structures are Trees and Graphs.
\ee

Another important division is between ``static'' and ``dynamic'' data structures.

\bd[Static Data Structures]
A \textbf{static data structure} is a data structure where the data in memory are fixed in size and allocated during
compile time.
\ed

In a static data structure the maximum size is needed to be known in advance, as memory cannot be reallocated at a
later point. Such data structures are easy to implement as computer memory is also sequential. Static data structures
have the advantage of fast access, but they are slower in insertion and deletion.

\be
An example of static data structure is an array.
\ee

\bd[Dynamic Data Structures]
A \textbf{dynamic data structure} is a data structure where the data in memory are allocated in real time.
\ed

A dynamic data structure refers to a collection of data in memory that has the flexibility to grow or shrink in size,
enabling a programmer to control exactly how much memory is utilised. Dynamic data structures have the advantage of
fast insertion and deletion, but they are much slower in access.

\be
An example of dynamic data structure is a linked-list.
\ee

Efficiency of data structures is always measured in terms of time and (memory) space. An ideal data structure would
be the one that takes the least possible time for all its operations and consumes the least memory space. \v

A data type is an attribute of data, as defined by the domain of values data can take, the programming language used,
or the operations that can be performed on the data.

\bd[Data Type]
A \textbf{data type} (or simply \textbf{type}) is an attribute of data which tells the compiler or interpreter how
the programmer intends to use the data. A data type constrains the values that a data might take and the operations
that can be done on the data.
\ed

In other words a data type provides a set of values from which an expression may take its values. \v

Most programming languages support basic data types called ``primitive data types''.

\bd[Primitive Data Type]
\textbf{Primitive data types} are types that are built-in or basic to a language implementation.
\ed

\be
Classic basic primitive types that are generally supported more or less directly by computer hardware and any
programming language are:
\bit
\item \textbf{Character}: is a unit of information that roughly corresponds to a symbol, such as in an alphabet, in the
written form of a natural language.
\item \textbf{Integer}: represents some range of mathematical integers that are commonly represented in a computer as
a group of bits. The size of the grouping varies so the set of integer sizes available varies between different types
of computers.
\item \textbf{Floating-point}: is an arithmetic unit using formulaic representation of real numbers as an approximation
to support a trade-off between range and precision.
\item \textbf{Fixed-point}: refers to a method of representing fractional numbers by storing a fixed number of digits of
their fractional part.
\item \textbf{Boolean}: is a data type that has one of two possible values (usually denoted ``true'' and ``false'')
which is intended to represent the two truth values of logic and Boolean algebra.
\item \textbf{Reference}: is a value that enables a program to indirectly access a particular data, such as a variable's
value or a record, in the computer's memory or in some other storage device.
\eit
\ee

Opposite to primitive data types, we also have the so called ``composite data types``.

\bd[Composite Data Type]
A \textbf{composite data type} is any data type which can be constructed in a program using the programming
language's primitive data types or other composite data types.
\ed